%% =====================================================================
%%  Corrigés — Unit 6 « A day in my life »
%% =====================================================================
\corrigesUnite{6}{A day in my life}

\begin{corrige}{1}
\begin{reponses}[2]
\item \rep{It's seven o'clock.}
\item \rep{It's eight fifteen.}
\item \rep{It's ten thirty.}
\item \rep{It's twelve forty-five.}
\item \rep{It's three twenty.}
\item \rep{It's six fifty.}
\end{reponses}
\pourquoi{en anglais, on n'utilise pas « 15 h » ou « 18 h » dans la conversation : 15 h 20 se dit
\en{three twenty} (et, si besoin, \en{three twenty p.m.}). On accepte aussi \en{It's a quarter past
eight}, \en{It's half past ten}, \en{It's a quarter to one}.}
\end{corrige}

\begin{corrige}{2}
\begin{reponses}[4]
\item \rep{works}
\item \rep{goes}
\item \rep{watches}
\item \rep{studies}
\item \rep{has}
\item \rep{plays}
\item \rep{finishes}
\item \rep{does}
\end{reponses}
\pourquoi{\en{play} garde son \en{y} car il suit une voyelle ; \en{study} le change en \en{ies}.
\en{have} → \en{has} est le seul verbe vraiment irrégulier au présent simple (avec \en{be}).}
\end{corrige}

\begin{corrige}{3}
\begin{reponses}[4]
\item \rep{gets up}
\item \rep{takes}
\item \rep{goes}
\item \rep{starts}
\item \rep{studies}
\item \rep{plays}
\item \rep{has}
\item \rep{finishes}
\end{reponses}
\retiens{avec \en{he} et \en{she}, tous les verbes prennent un \en{-s} : on le vérifie un par un en
relisant son texte.}
\end{corrige}

\begin{corrige}{4}
\begin{reponses}[3]
\item \rep{works}
\item \rep{take}
\item \rep{drink}
\item \rep{walks}
\item \rep{have}
\item \rep{watches}
\end{reponses}
\pourquoi{\en{Jack and Emma} = \en{they} : pas de \en{-s}. \en{my mother}, \en{Léa}, \en{Mrs.\ Lopez}
= \en{she} : \en{-s}.}
\end{corrige}

\begin{corrige}{5}
\begin{reponses}[3]
\item \rep{Tuesday}
\item \rep{Thursday}
\item \rep{Saturday}
\item \rep{Sunday}
\item \rep{Wednesday}
\item \rep{on Mondays}
\end{reponses}
\piege{oublier la majuscule des jours, ou écrire « Wendsday » : on écrit \en{Wednesday} (et l'on dit
\phon{ˈwenzdeɪ}).}
\end{corrige}

\begin{corrige}{6}
\begin{reponses}[6][french]
\item \rep{b}
\item \rep{e}
\item \rep{d}
\item \rep{a}
\item \rep{f}
\item \rep{c}
\end{reponses}
\pourquoi{l'ordre suit les heures : 7 h 30, 8 h, 9 h, 13 h, 18 h, 23 h.}
\end{corrige}

\begin{corrige}{7}
\begin{reponses}[3]
\item \rep{/s/}
\item \rep{/z/}
\item \rep{/ɪz/}
\item \rep{/s/}
\item \rep{/z/}
\item \rep{/ɪz/}
\item \rep{/s/}
\item \rep{/z/}
\item \rep{/ɪz/}
\end{reponses}
\pourquoi{c'est la même règle que pour le pluriel des noms. Texte de la piste page~\pageref{trans-6.10}.}
\end{corrige}

\begin{corrige}{8}
\begin{reponses}[4]
\item \rep{8:00}
\item \rep{8:40}
\item \rep{9:15}
\item \rep{1:30}
\item \rep{6:00}
\item \rep{7:00}
\item \rep{11:30}
\end{reponses}
\pourquoi{\en{eight forty} = 8 h 40 ; \en{one thirty} = 13 h 30 ; \en{eleven thirty} = 23 h 30.
Emma reconnaît : \en{I'm sometimes late!} (\en{sometimes} = parfois).}
\end{corrige}

\begin{corrige}{9}
\begin{reponses}[4]
\item \rep{gets}
\item \rep{goes}
\item \rep{has}
\item \rep{takes}
\item \rep{starts}
\item \rep{has}
\item \rep{watches}
\item \rep{goes}
\end{reponses}
\retiens{dans un récit à la 3\textsuperscript{e} personne, chaque verbe porte un \en{-s} : tendez
l'oreille, on l'entend toujours.}
\end{corrige}

\begin{corrige}{10}
\begin{reponses}[1]
\item \rep{What time is it? — It's eight thirty.} \textit{ou} \rep{It's half past eight.}
\item \rep{I get up at seven every day.}
\item \rep{Emma takes the subway.}
\item \rep{He finishes work at six.}
\item \rep{On Saturdays, I sleep all day.}
\item \rep{My sister studies music.}
\item \rep{We have lunch at noon.}
\item \rep{Jack watches TV in the evening.}
\end{reponses}
\piege{« Je me lève » : pas de \en{be}, et pas de pronom réfléchi. \en{get up} suffit.}
\end{corrige}

\begin{corrige}{11}
Réponse modèle :
\begin{anglais}
\begin{tcolorbox}[otBase,colback=otGreenTint,colframe=otGreenTint,fontupper=\fontsize{10.4}{14}\selectfont]
I get up at seven. I take a shower and I have breakfast at seven thirty. I take the subway at
eight fifteen and I start work at nine. I have lunch at one. I finish work at six and I go home.
In the evening, I cook dinner and I watch TV. I go to bed at eleven.
\end{tcolorbox}
\end{anglais}
\end{corrige}

\begin{corrige}{12}
Réponses modèles :
\begin{reponses}[1]
\item \rep{It's nine fifteen.}
\item \rep{I get up at seven. I have breakfast and I take the subway.}
\item \rep{I work. I have lunch at one.}
\item \rep{I cook dinner and I watch TV.}
\item \rep{Yes! It's a long day!}
\end{reponses}
\end{corrige}

\begin{corrigetest}
\begin{reponses}[5][french]
\item \rep{b}
\item \rep{a}
\item \rep{c}
\item \rep{a}
\item \rep{b}
\item \rep{c}
\item \rep{b}
\item \rep{a}
\item \rep{c}
\item \rep{b}
\end{reponses}
\pourquoi{question 1 : on demande l'heure avec \en{time}, jamais avec \en{hour}. Question 4 :
\en{they} n'a pas de \en{-s}. Question 10 : \en{p.m.} = l'après-midi et le soir.}
\end{corrigetest}
